\section{Results: what the engine returns}
\label{sec:results}

Every number here is engine behaviour under evidence tier 5 inputs. All sixteen regulatory gates
carry status \texttt{UNCERTAIN} with no reviewer, every run returns eligibility
\texttt{NO\_CONCLUSION} for every strategy, and nothing below measures any plant, product or
supply chain. The frozen requirement is a tail statement, not a mean: a design meets it when
$P(\text{run-level fill rate} \ge 0.99) \ge 0.90$ over a five-year measured window after a
365-day warm-up. Both numbers were pre-specified on 2026-09-01 as revision R001, before any
simulation result existed. Among designs that meet it, lower annual total cost is preferred.
Cost never buys a relaxation of the requirement. Because the requirement is a Monte Carlo
proportion, the binomial standard error at the requirement governs what a verdict is worth: the
optimization and the ablations run at 60 paired runs (0.039), the dominance and reversal maps at
40 (0.047), the threshold bisections at 20 (0.067), the simulation and the re-evaluation at 100
(0.030). All postdate revisions R009 to R011, so earlier runs are superseded and none appears
here.

One structural fact separates the two products in almost every table below. At the illustrative
inputs the norepinephrine network's status-quo saleable capacity runs at 0.778 of demand-implied
throughput and the sodium bicarbonate network's at 1.245. The bicarbonate case is capacity-short
by construction, and a network whose saleable capacity is below demand cannot be stocked out of
the deficit over a long enough window. That is a property of the illustrative demand and capacity
numbers, not an observation about the products.

\begin{table}[htbp]
\centering
\caption{Specification outputs at each architecture's optimized design, 60 paired runs,
\texttt{opt\_20260906T043357Z}. Days avoided is the paired reduction in shortage days per year
against S0 on the same seeds; incremental cost is thousands of USD per shortage-day avoided per
year and reads ``none'' where no days are avoided. Response times are medians across episodes.
Six architectures are omitted for space and appear in Table~\ref{tab:feasible} where they
qualify. All inputs tier 5 illustrative.}
\label{tab:outputs}
\scriptsize
\setlength{\tabcolsep}{4pt}
\begin{tabular}{@{}lrrrrrrrrr@{}}
\toprule
 & USD/ & Annual & Mean & $P$(meet & Short. & Days & Incr. & First & Stable \\
Id & unit & USD M & fill & dem.) & d/yr & avoid. & kUSD/d & resp. d & recov. d \\
\midrule
\multicolumn{10}{@{}l}{\textit{Norepinephrine bitartrate 1 mg/mL, 4 mL vial (resolved comparator)}}\\
S0 & 13.82 & 12.91 & 0.9949 & 0.917 & 7.5 & +0.0 & n/a & 3.7 & 4.0 \\
S1 & 13.82 & 12.91 & 0.9949 & 0.917 & 7.5 & +0.0 & none & 3.7 & 4.0 \\
S2 & 19.22 & 18.03 & 0.9986 & 0.950 & 1.5 & +6.0 & 848 & 1.1 & 8.4 \\
S3 & 18.70 & 17.48 & 0.9952 & 0.900 & 9.4 & $-1.9$ & none & 1.0 & 2.5 \\
S4 & 21.53 & 20.18 & 0.9975 & 0.950 & 2.1 & +5.4 & 1347 & 0.7 & 10.3 \\
S5 & 21.39 & 20.04 & 0.9975 & 0.950 & 5.3 & +2.2 & 3266 & 8.3 & 12.8 \\
S6 & 22.45 & 21.03 & 0.9975 & 0.950 & 5.3 & +2.2 & 3721 & 8.3 & 12.8 \\
S7 & 16.93 & 15.83 & 0.9956 & 0.933 & 6.6 & +1.0 & 3021 & 6.7 & 4.3 \\
S9 & 24.85 & 23.28 & 0.9972 & 0.917 & 2.0 & +5.6 & 1864 & 2.6 & 12.7 \\
S11 & 16.22 & 15.15 & 0.9949 & 0.933 & 3.9 & +3.7 & 611 & 2.8 & 19.3 \\
S15 & 13.76 & 12.87 & 0.9958 & 0.917 & 9.1 & $-1.6$ & none & 2.1 & 3.5 \\
S16 & 17.07 & 16.00 & 0.9981 & 0.950 & 1.3 & +6.2 & 497 & 1.9 & 6.4 \\
S17 & 19.64 & 18.38 & 0.9967 & 0.950 & 10.7 & $-3.1$ & none & 6.0 & 7.5 \\
S19 & 17.02 & 15.88 & 0.9939 & 0.917 & 3.8 & +3.7 & 806 & 2.2 & 18.5 \\
\midrule
\multicolumn{10}{@{}l}{\textit{Sodium bicarbonate 8.4\%, 50 mL vial (prolonged current shortage)}}\\
S0 & 11.69 & 13.41 & 0.9200 & 0.017 & 64.2 & +0.0 & n/a & 14.8 & 0.0 \\
S1 & 11.69 & 13.41 & 0.9200 & 0.017 & 64.2 & +0.0 & none & 14.8 & 0.0 \\
S2 & 19.26 & 24.00 & 0.9949 & 0.850 & 15.3 & +48.9 & 217 & 3.9 & 19.6 \\
S3 & 14.82 & 18.49 & 0.9963 & 0.917 & 3.5 & +60.7 & 84 & 1.3 & 10.7 \\
S4 & 19.23 & 24.04 & 0.9984 & 0.950 & 1.2 & +63.0 & 169 & 0.9 & 7.6 \\
S5 & 25.63 & 31.98 & 0.9961 & 0.900 & 3.6 & +60.6 & 306 & 3.8 & 13.2 \\
S6 & 27.02 & 33.70 & 0.9961 & 0.900 & 3.6 & +60.6 & 335 & 3.8 & 13.2 \\
S7 & 13.77 & 16.56 & 0.9632 & 0.317 & 32.8 & +31.4 & 100 & 8.0 & 0.3 \\
S9 & 19.68 & 24.63 & 0.9993 & 0.967 & 1.0 & +63.2 & 177 & 0.8 & 6.0 \\
S11 & 12.89 & 16.11 & 0.9982 & 0.950 & 1.9 & +62.3 & 43 & 0.3 & 4.8 \\
S15 & 11.77 & 13.38 & 0.9122 & 0.017 & 69.3 & $-5.1$ & none & 14.9 & 0.0 \\
S16 & 13.40 & 16.72 & 0.9961 & 0.917 & 3.8 & +60.4 & 55 & 1.4 & 24.2 \\
S17 & 15.57 & 19.41 & 0.9958 & 0.917 & 15.4 & +48.8 & 123 & 4.4 & 8.4 \\
S19 & 14.77 & 16.35 & 0.8881 & 0.000 & 99.9 & $-35.7$ & none & 13.3 & 271.9 \\
\bottomrule
\end{tabular}
\end{table}

Four things follow from Table~\ref{tab:outputs}, each a statement about the model. First, on
norepinephrine the status quo at its own optimized inventory policy already returns mean fill
0.9949 and $P(\text{meet}) = 0.917$ at 13.82 USD per unit, where the configured default design
in the 100-run simulation returns 0.9474 and 0.100 at 14.28 USD per unit. The two come from runs
at different resolutions, so only the direction should be read, but the direction is large: on
the product whose plant is not capacity-short, most of the achievable service comes from
inventory policy at the incumbent plant, at roughly 0.2 million USD per year, with no new plant,
no second source and no new site. Second, once that baseline exists, further service is
expensive: the cheapest architecture that avoids any shortage days on norepinephrine does so at
497 thousand USD per shortage-day avoided per year, and distributed regional nodes at 3.27
million, where on sodium bicarbonate, whose baseline loses 64 shortage days a year, the same two
figures are 43 thousand and 306 thousand. The specification's headline efficiency output differs
by nearly two orders of magnitude between the products, and the difference is whether the
incumbent network has enough capacity. Third, the release-assurance variant costs money and
returns nothing measurable: S6 matches S5 to four decimal places on both products while costing
0.99 and 1.73 million USD a year more, and part of that identity is defect MD-18, so the cost
difference is real inside the ledger while the service equality is partly an artifact. Fourth,
response time is not the binding output: medians to first response run from 0.3 to 14.9 days and
to stable recovery from 0.0 to 271.9, and recovery is slowest where fill is already high, because
an architecture that rarely stocks out records few episodes and clears each slowly relative to
its own small deficit.

\begin{table}[htbp]
\centering
\caption{The verified feasible set at 100 paired runs, from
\texttt{results/design\_space/reeval\_post\_R009.json}. A cell is borderline when its estimated
$P(\text{meet})$ lies within one binomial standard error (0.030) of the 0.90 requirement.}
\label{tab:feasible}
\footnotesize
\setlength{\tabcolsep}{5pt}
\begin{tabular}{@{}llrrrrl@{}}
\toprule
Product & Id & Annual USD M & Mean fill & $P$(meet) & SE & Borderline \\
\midrule
Norepinephrine & S11 & 15.2 & 0.9952 & 0.900 & 0.030 & yes \\
Norepinephrine & S19 & 15.9 & 0.9943 & 0.920 & 0.027 & yes \\
Norepinephrine & S2 & 18.0 & 0.9989 & 0.960 & 0.020 & no \\
Norepinephrine & S17 & 18.4 & 0.9975 & 1.000 & 0.000 & no \\
Norepinephrine & S13 & 18.8 & 0.9964 & 0.900 & 0.030 & yes \\
Norepinephrine & S5 & 20.0 & 0.9963 & 0.940 & 0.024 & no \\
Norepinephrine & S6 & 21.0 & 0.9963 & 0.940 & 0.024 & no \\
Norepinephrine & S8 & 22.5 & 0.9964 & 0.910 & 0.029 & yes \\
Norepinephrine & S9 & 23.3 & 0.9975 & 0.920 & 0.027 & yes \\
Norepinephrine & S14 & 23.9 & 0.9964 & 0.920 & 0.027 & yes \\
Bicarbonate & S17 & 19.4 & 0.9960 & 0.900 & 0.030 & yes \\
Bicarbonate & S12 & 21.5 & 0.9967 & 0.920 & 0.027 & yes \\
Bicarbonate & S18 & 24.2 & 0.9977 & 0.930 & 0.026 & no \\
Bicarbonate & S9 & 24.6 & 0.9987 & 0.940 & 0.024 & no \\
\bottomrule
\end{tabular}
\end{table}

Table~\ref{tab:feasible} gives the fourteen of forty product-strategy cells meeting the
requirement when each cell's own incumbent design is re-evaluated at 100 paired runs without
re-searching, which separates a change of classification caused by a model fix from one caused by
a small final-run count. Eight of the fourteen are borderline, and eight further cells the run
calls infeasible also sit within one standard error, sixteen of forty in total. What that means
is narrow: at 100 paired runs a cell estimated at 0.900 and one at 0.870 cannot be told apart, so
their opposite labels carry no information, and only six cells across both products are separated
from the requirement by more than one standard error. Two architectures meet it on both products,
rotating prequalified campaigns (S17) and dual finished-dose sourcing with a
key-starting-material tier (S9), and S17 returns the only tail probability of 1.000 in the set.
The table is a floor rather than a boundary, because it scores incumbents of an earlier and
smaller search. The later search over the same space finds a feasible design for fifteen of the
twenty-six cells this re-evaluation calls infeasible, including sodium bicarbonate S11 at 0.950
against 0.630 and norepinephrine S0 itself. The disagreement runs one way only, which is the
signature of a search that has not converged rather than of a model that has changed. The honest
reading is that this package can identify architectures that plainly fail and cannot yet rank the
ones that do not.

No price is assumed anywhere. The break-even price is derived from each run's own cost ledger and
delivered volume, and the break-even volume is the committed volume covering fixed and committed
cost at that price; no reimbursement figure enters the cost side or stands in for a price. Across
all forty cells the fixed share of annual cost runs from 0.797 to 0.937, so every architecture
here is a fixed-cost business in the model and its break-even volume is a large fraction of its
own capacity: 0.64 to 0.94 overall, and 0.80 or above for every architecture meeting the
requirement on either product. On norepinephrine the break-even price is 14.26 USD per unit for
the status quo, 16.63 for bright-stock campaign offtake, 20.55 for rotating campaigns, 46.65 for
distributed regional nodes and 49.10 once the release layer is added; on sodium bicarbonate the
same five are 14.44, 13.28, 15.72, 40.59 and 42.67. These are model cost recoveries under
illustrative inputs, not prices, not offers, and not evidence that any buyer would pay them.

\section{Sensitivity: what determines the answer}
\label{sec:sensitivity}

Three mechanisms decide whether an architecture in this battery meets the requirement, and they
are not the three the original thesis rests on. In order of measured effect they are capacity
adequacy, the regional inventory review rule, and commissioning time, and capacity adequacy
dominates the other two by roughly an order of magnitude on every service metric. Release time,
the mechanism a release-assurance layer would address, is worth two of the sixteen days in the
release path and moves mean fill by less than 0.002 anywhere in the battery. Capital cost and
fixed quality and labour cost change feasibility in none of the forty cells at any point of their
declared ranges, and in the reversal map they do not change even the ranking.

\begin{table}[htbp]
\centering
\caption{One-way threshold census over 40 product-strategy cells per input, 480 bisections at 20
paired runs each, from \texttt{ds\_post\_R009/thresholds.csv}. The last two columns count cells
whose classification never changes anywhere in the declared range.}
\label{tab:census}
\footnotesize
\setlength{\tabcolsep}{5pt}
\begin{tabular}{@{}llrrr@{}}
\toprule
Input & Declared range & Thresh. & Feasible & Infeas. \\
      &                & inside  & throughout & throughout \\
\midrule
Annual demand & 0.5M to 1.8M; 0.6M to 2.4M & 40 & 0 & 0 \\
Batches per site-year & 30 to 70 & 39 & 1 & 0 \\
Common-cause events per year & 0.02 to 0.3 & 31 & 4 & 5 \\
Changeover days & 1 to 10 & 25 & 3 & 12 \\
Shelf life, months & 12 to 30; 12 to 36 & 19 & 9 & 12 \\
Raw-material lead time, days & 30 to 240 & 19 & 11 & 10 \\
Supplier concentration & 0.05 to 0.6 & 17 & 14 & 9 \\
Release time, sterility, days & 14 to 18 & 6 & 22 & 12 \\
Node commissioning, days & 365 to 1095 & 4 & 24 & 12 \\
Contract activation latency, days & 7 to 60 & 2 & 26 & 12 \\
Capital per node, USD & 15M to 100M & 0 & 28 & 12 \\
Fixed quality and labour, USD/site-yr & 1.5M to 6M & 0 & 28 & 12 \\
\bottomrule
\end{tabular}
\end{table}

Two rows of Table~\ref{tab:census} carry the result. Annual demand changes feasibility in all
forty cells and batches per site-year in thirty-nine of forty; both are statements about capacity
adequacy, the first varying the demand a fixed network must serve and the second how much it can
make. No other input changes feasibility in even four fifths of the cells, and capital per node
and fixed quality and labour per site-year change it in zero cells across a four-fold and a
six-and-a-half-fold range, so cost is not a feasibility lever in this model. Two cautions belong
with the census: the bisections run at 20 paired runs, where the standard error is 0.067, so a
threshold is an interval and not a point; and twelve cells sit in the last column for most inputs
and can never show a threshold, so the informative denominator for the weaker inputs is 28 rather
than 40. As multiples of the base demand each cell was designed for, the demand thresholds run
from 0.71 to 1.17, median 1.075 on norepinephrine and 0.99 on sodium bicarbonate. No architecture
buys more than about seventeen per cent of demand headroom over its own design point, and the
largest headroom in the battery is the status quo's own. Five norepinephrine cells lose
feasibility below the demand they were designed for, and on sodium bicarbonate fourteen of twenty
do. These are sized designs rather than resilient ones, sitting near a capacity cliff that the
choice between them moves by a few per cent of demand.

Two of the seven sensitivity axes the specification names are covered only indirectly, which is
stated rather than papered over. Uptime and yield enter the model only through saleable capacity
in Eq.~\ref{eq:capacity}, so a proportional change in either is the same lever as the same
proportional change in batches per site-year, the row that already moves 39 of 40 cells; neither
was bisected separately in the threshold grid. Routine demand volatility likewise was not
bisected: it appears in the ablation below as the demand-shock and correlated-regional-shock
mechanisms, both of which rank low, while the demand \emph{level} was swept and dominates. The
other five axes, capacity utilization, release time, quality and fixed cost, and raw-material lead
time, are swept directly in Table~\ref{tab:census} or ablated in Table~\ref{tab:ablation}.

\begin{table}[htbp]
\centering
\caption{Leave-one-out attribution over 40 cells, 60 paired runs, from
\texttt{abl\_post\_R009/attribution.csv}. Positive deltas mean removing the mechanism improves
service; cost is the mean change in annual total cost, thousands of USD. Five mechanisms with
mean $\Delta$fill below 0.0003 are omitted, and the three cost mechanisms, which move every
service metric by exactly zero, are reported as one row.}
\label{tab:ablation}
\scriptsize
\setlength{\tabcolsep}{4pt}
\begin{tabular}{@{}llrrrrrr@{}}
\toprule
Mechanism removed & Family & Mean & Max & Mean & Max & Short. & Cost \\
                  &        & $\Delta$fill & $\Delta$fill & $\Delta P$ & $\Delta P$ & d/yr & kUSD \\
\midrule
Capacity shortfall & capacity & +0.01128 & +0.11066 & +0.150 & +0.967 & $-10.2$ & +379 \\
Raw-material buffer & supply & +0.00535 & +0.03508 & +0.049 & +0.167 & $-4.2$ & +535 \\
Common-cause events & dependence & +0.00522 & +0.04325 & +0.052 & +0.233 & $-4.6$ & +57 \\
Supplier disruption & supply & +0.00393 & +0.02137 & +0.031 & +0.083 & $-2.9$ & +33 \\
Idiosyncratic site failure & capacity & +0.00371 & +0.02342 & +0.041 & +0.217 & $-3.0$ & +41 \\
Backorder window, 7 to 30 d & structure & +0.00312 & +0.01293 & +0.039 & +0.117 & +2.1 & +4 \\
Deviations, rejections, yield var. & quality & +0.00148 & +0.00841 & +0.015 & +0.100 & $-1.3$ & $-13$ \\
Regional review rule & inventory & +0.00143 & +0.00887 & +0.032 & +0.233 & $-3.5$ & +141 \\
Demand shocks & demand & +0.00127 & +0.00860 & +0.004 & +0.083 & $-1.1$ & $-9$ \\
Sterility incubation, 14 days & release & +0.00099 & +0.00323 & +0.013 & +0.067 & $-1.4$ & +17 \\
Correlated regional shocks & demand & +0.00084 & +0.00493 & +0.004 & +0.033 & $-0.7$ & $-4$ \\
Commissioning delay & time & +0.00054 & +0.00495 & +0.015 & +0.150 & $-0.8$ & +285 \\
503B availability & regulatory & +0.00034 & +0.01200 & +0.004 & +0.133 & $-0.3$ & +5 \\
Chemical release queue, 2 days & release & +0.00031 & +0.00149 & +0.000 & +0.033 & $-0.5$ & +6 \\
The three cost mechanisms & cost & 0.00000 & 0.00000 & 0.000 & 0.000 & 0.0 & $-5431$ \\
Component lead time & supply & $-0.00101$ & +0.00053 & $-0.005$ & +0.017 & +0.6 & $-15$ \\
Ten-year horizon & structure & $-0.01504$ & +0.00276 & $-0.051$ & +0.083 & +12.7 & +375 \\
\bottomrule
\end{tabular}
\end{table}

Capacity shortfall is more than twice the next mechanism on mean fill and three times it on the
tail probability, and its largest single-cell effect, +0.111 fill and +0.967 on $P(\text{meet})$
at sodium bicarbonate S19, is an order of magnitude larger than any other mechanism's largest.
Two rows are warnings rather than findings. Component lead time carries a negative sign, meaning
a shorter lead slightly worsens service; revision R011 fixed the policy half of that defect and a
residual remains, because under a forced outage a deferred order lands when the supplier resumes
whatever its nominal lead, so a longer lead carries a larger protective pipeline. That is a
property of the disruption model and is recorded rather than argued away. The ten-year horizon
row is a measurement convention, not a mechanism.

\begin{table}[htbp]
\centering
\caption{Interaction of the three deciding mechanisms, 60 paired runs, averaged over 40 cells,
from \texttt{abl\_post\_R009\_pairs/}. The additive prediction adds the single-mechanism deltas
to the base; the last two columns are observed minus predicted.}
\label{tab:interaction}
\footnotesize
\setlength{\tabcolsep}{4pt}
\begin{tabular}{@{}lrrrrrrr@{}}
\toprule
Mechanisms removed & Mean & $P$(meet) & Short. & Annual & Additive & Fill & $P$ \\
                   & fill &           & d/yr   & USD M  & fill     & inter. & inter. \\
\midrule
None (base) & 0.98666 & 0.8058 & 13.60 & 19.48 & . & . & . \\
Capacity adequacy only & 0.99794 & 0.9558 & 3.37 & 19.86 & 0.99794 & 0.00000 & 0.0000 \\
Regional review rule only & 0.98809 & 0.8379 & 10.10 & 19.62 & 0.98809 & 0.00000 & 0.0000 \\
Commissioning only & 0.98720 & 0.8208 & 12.80 & 19.76 & 0.98720 & 0.00000 & 0.0000 \\
Capacity plus review rule & 0.99906 & 0.9796 & 0.69 & 20.03 & 0.99937 & $-0.00031$ & $-0.0083$ \\
Capacity plus commissioning & 0.99811 & 0.9617 & 3.10 & 20.14 & 0.99847 & $-0.00036$ & $-0.0092$ \\
Review rule plus commissioning & 0.98842 & 0.8467 & 9.80 & 19.91 & 0.98863 & $-0.00021$ & $-0.0062$ \\
All three & 0.99913 & 0.9817 & 0.62 & 20.33 & 0.99990 & $-0.00078$ & $-0.0212$ \\
\bottomrule
\end{tabular}
\end{table}

Ranking mechanisms one at a time does not show whether they substitute. Removing all three lifts
the average tail probability from 0.806 to 0.982 and cuts average shortage days from 13.6 to 0.6,
at an average cost increase of 0.85 million USD per year. The interactions are all negative and
all small, so the three substitute weakly rather than reinforce: fixing capacity first removes
most of the value of fixing the other two, and capacity alone recovers 0.0113 of the 0.0125 fill
and 0.150 of the 0.176 tail probability the full triple delivers. That ordering is the practical
content of the sensitivity analysis, and it exposes a gap: the raw-material buffer ranks second in
Table~\ref{tab:ablation} and is a design variable in every search space, but it was left out of
the interaction study, so its substitution behaviour with capacity is unmeasured.

\begin{table}[htbp]
\centering
\caption{Decision-reversal map, from \texttt{ds\_post\_R009/reversal.json}, 40 paired runs per
cell. Each entry is the preferred strategy with the number of feasible strategies of 20 in
brackets; demand is a multiple of the base for that product. Entries were identical at all three
levels of fixed quality and labour cost per site-year (1.5M, 3M and 6M USD), so that axis is
collapsed.}
\label{tab:reversal}
\footnotesize
\setlength{\tabcolsep}{8pt}
\begin{tabular}{@{}llll@{}}
\toprule
Product and demand & Common-cause 0.02 & 0.10 & 0.30 \\
\midrule
Norepinephrine, 0.56 & S15 (19) & S15 (19) & S15 (18) \\
Norepinephrine, 1.00 & S15 (20) & S15 (20) & S17 (1) \\
Norepinephrine, 2.00 & none (0) & none (0) & none (0) \\
Bicarbonate, 0.50 & S0 (20) & S15 (20) & S15 (19) \\
Bicarbonate, 1.00 & S11 (15) & S11 (15) & none (0) \\
Bicarbonate, 2.00 & none (0) & none (0) & none (0) \\
\bottomrule
\end{tabular}
\end{table}

Table~\ref{tab:reversal} contains three results. A four-fold change in the fixed quality and
labour cost of a site reverses no decision anywhere: in a study about whether to build
manufacturing capacity, the cost of running a quality unit does not decide the answer and the
amount of capacity does. At twice the design demand nothing in the battery is feasible on either
product, which is the demand cliff seen from the other side. And dependence between sites is the
one hazard that can collapse the feasible set on its own: at base demand and a common-cause rate
of 0.30 per year, norepinephrine drops from twenty feasible strategies to one and sodium
bicarbonate from fifteen to none. Correlated failure rather than site count is what an
architecture has to defeat, and adding sites that still share an API source, a container supplier,
a geography or a quality unit does not defeat it. Most rows prefer S15, which builds no new plant,
because it is the cheapest design meeting the requirement, which is the frozen decision rule
working as specified. Two qualifications stop that being a recommendation: the map runs at 40
paired runs, where the standard error is 0.047, so its feasible sets are coarse; and S15 is a
product-selection construct, a presentation whose capacity is adequate and whose buffer is
payable, so its position at the top says which product to pick and not which network to build.

On release time specifically, a chemical release-assurance layer can only reduce the assay
component of the sixteen-day path, and the engine's ablation for it zeroes assay, endotoxin,
environmental monitoring and quality review together, removing 2 of 16 days for +0.00031 mean
fill and +0.00149 in the single cell where it does most. Removing the entire 14-day sterility
incubation, which no release-assurance layer can do and which only a change of sterilization
route with parametric release could reach, moves mean fill by +0.00099 and at most +0.00323. A
perfect quality system with no deviations, rejections or yield variance moves it by +0.00148 and
at most +0.00841, the largest figure this package has recorded for any release-side or
quality-side ablation in any run vintage. On mean fill, capacity is worth eight times the quality
ablation and eleven times the sterility one, and the one-way analysis agrees: the sterility
component changes feasibility in 6 of 40 cells against 40 of 40 for demand. In the dominance
analysis S6 is dominated by S5 on both products. Defect MD-18 means part of that identity is an
artifact, but no run in any vintage has shown the release layer paying for itself.

Four limits on this analysis are load-bearing and none is closed. The inputs are illustrative, so
it says what the model is sensitive to and not where the real world sits inside those ranges; the
most consequential unknown is the one the analysis calls decisive, the true capacity headroom of
the incumbent networks for these two presentations, and the package holds no evidence for it above
tier 5. The resolution is coarse and inconsistent, at 20, 40, 60 and 100 paired runs with standard
errors of 0.067, 0.047, 0.039 and 0.030, and fifteen of forty cells change classification between
the two most recent optimizations, so individual cell verdicts are provisional throughout. Several
model defects bearing on these numbers remain open and are listed in Table~\ref{tab:defects}. And
nothing here has been reviewed by anyone.
